Welcome to the 11th Social Media Mining for Health Research and Applications and Health Real-World Data (\#SMM4H-HeaRD) Workshop and Shared Tasks, co-located with the 64th Annual Meeting of the Association for Computational Linguistics (ACL 2026). This year, \#SMM4H-HeaRD 2026 will be held online, continuing our tradition of connecting data mining researchers who focus on leveraging social media and real-world health data for health informatics.

\par\vspace{1em}
This year's edition emphasizes privacy-preserving data sharing and real-world evaluation, particularly in the context of health-related text from sources such as social media, electronic health records, and biomedical literature. In alignment with the ACL 2026 Theme Track on the Explainability of NLP Models, we also invited work that improves the transparency, interpretability, and trustworthiness of models applied in high-stakes clinical and public health contexts.

\par\vspace{1em}
For \#SMM4H-HeaRD 2026, we organized eight shared tasks spanning diverse domains and data modalities, including detection of adverse drug events in multilingual social media posts, detection of insomnia in clinical notes, estimating flu vaccine effectiveness from social media, generation of structured medical notes from dialogues, detection of patient metadata in SARS-CoV-2 sequencing articles, predicting TNM staging from pathology reports, extraction of social and clinical impacts of substance use, and multilingual clinical entity annotation projection and extraction. Each submission underwent a rigorous review process by program committee members chosen for their expertise.

\par\vspace{1em}
We hope you find the workshop papers insightful and inspiring. We extend our gratitude to the program committee, additional reviewers, ACL 2026 organizers, annotators of the shared task data, and everyone who submitted a paper or participated in the shared tasks. \#SMM4H-HeaRD 2026 would not have been possible without their contributions.

\par\vspace{2em}
\#SMM4H-HeaRD 2026 Workshop Chairs and Shared Task Committee
